% ===========================================================================
%  Ficha XGodel — Divisores · Primero de Secundaria · Aritmética
%  Compilar:  xelatex ficha-divisores.tex   (dos pasadas)
% ===========================================================================
\documentclass{xgodel-ficha}

\marca{XGodel}
\sitio{www.XGodel.com}
\grado{Primero de Secundaria}
\curso{Aritmética}
\titulo{Divisores}

\begin{document}
\portada

\begin{ficha}

% --------------------------- COLUMNA IZQUIERDA ---------------------------
\begin{pensamiento}
Si puedes repartir $12$ caramelos en grupos exactos, ¡ya sabes buscar
divisores! Hoy aprenderás a encontrarlos todos sin equivocarte.
\end{pensamiento}

\concepto{Divisor de un número}
Un número $d$ es \textbf{divisor} de $N$ si la división $N \div d$ es
\textbf{exacta}, es decir, el residuo es cero.
\lineas{2}

\formula{$N \div d = k$ \quad y \quad residuo $= 0$}

\concepto{Ejemplo resuelto}
Hallemos todos los divisores de $12$. Probamos desde el $1$:
\par\smallskip
$12 \div 1 = 12$ \quad $12 \div 2 = 6$ \quad $12 \div 3 = 4$
\par\smallskip
$12 \div 4 = 3$ \quad $12 \div 6 = 2$ \quad $12 \div 12 = 1$
\par\smallskip
En cambio $12 \div 5$, $12 \div 7$, $12 \div 8$, $12 \div 9$, $12 \div 10$
y $12 \div 11$ dejan residuo. Entonces:
\formula{$D_{12} = \{1,\, 2,\, 3,\, 4,\, 6,\, 12\}$}
El $12$ tiene $6$ divisores. Nota que el $1$ y el mismo número
\textbf{siempre} son divisores.

\begin{globo}
Ahora, ¿cómo saber si un número es divisor \textbf{sin hacer la división}?
Usamos los \textbf{criterios de divisibilidad}.
\end{globo}

\concepto{Criterios de divisibilidad}

\begin{apartados}

\item \rotulo{Divisibilidad por 2}
  El número termina en cifra par: $0$, $2$, $4$, $6$ u $8$.
  \par\smallskip
  Ejemplos: $34$ \textbf{sí} (termina en $4$); \; $57$ \textbf{no}.

\item \rotulo{Divisibilidad por 3}
  La suma de sus cifras es múltiplo de $3$.
  \par\smallskip
  Ejemplo: $231 \rightarrow 2 + 3 + 1 = 6$, y $6$ es múltiplo de $3$;
  entonces $231$ \textbf{sí}.
  \par\smallskip
  Ejemplo: $124 \rightarrow 1 + 2 + 4 = 7$; entonces $124$ \textbf{no}.

\item \rotulo{Divisibilidad por 5}
  El número termina en $0$ o en $5$.
  \par\smallskip
  Ejemplos: $145$ \textbf{sí}; \; $62$ \textbf{no}.

\end{apartados}

% ======================= EJERCICIOS DE APLICACIÓN =========================
\bloque{Ejercicios de Aplicación}

\begin{ejercicios}

\item Escribe todos los divisores de $15$:
  \par\smallskip
  $D_{15} = \{$ \rule{34mm}{0.4pt} $\}$

\item ¿Cuántos divisores tiene el número $24$?
  \alts{4}{6}{7}{8}{10}

\item Halla la suma de todos los divisores de $8$.
  \alts{12}{14}{15}{16}{18}

\item ¿Cuál de los siguientes números es divisible entre $3$?
  \alts{124}{235}{502}{411}{617}

\item Halla el mayor divisor de $20$ que sea menor que $20$.
  \alts{2}{4}{5}{10}{20}

\item ¿Cuántos números del $1$ al $20$ son divisibles entre $5$?
  \alts{3}{4}{5}{6}{10}

\end{ejercicios}

\end{ficha}
\end{document}
